% TOP_PROBLEM: {"problemId": "problem.parshin-conjecture", "canonicalTitle": "Parshin Conjecture on Higher K-Theory over Finite Fields", "releaseRank": 222, "primaryDomain": "number_theory_arithmetic_geometry", "primaryDomainLabel": "Number theory & arithmetic geometry", "displayStatus": "open", "releaseStatus": "open"}
% !TeX program = lualatex
% Definition and sources only; this file is not an LLM solution attempt.
\documentclass[11pt]{article}
\usepackage[margin=25mm]{geometry}
\usepackage{fontspec}
\setmainfont{Latin Modern Roman}
\usepackage{amsmath,amssymb,mathtools}
\usepackage{unicode-math}
\setmathfont{Latin Modern Math}
\usepackage{xurl,hyperref}
\hypersetup{colorlinks=true,urlcolor=blue}
\setcounter{secnumdepth}{0}
\setlength{\parindent}{0pt}
\setlength{\parskip}{5pt}
\setlength{\emergencystretch}{3em}
\begin{document}
\section*{\texorpdfstring{Parshin Conjecture on Higher K-Theory over Finite Fields}{Parshin Conjecture on Higher K-Theory over Finite Fields}}\label{222}
\subsection{Definitions and mathematical statement}
For a smooth projective variety $X$ over a finite field ${\mathbb{F}}_q$, let $K_n(X)$ denote its higher algebraic $K$-group. An abelian group is torsion when every element is killed by some positive integer, with the integer allowed to depend on the element. Parshin's assertion is
\[
 K_n(X)\otimes_{{\mathbb{Z}}}{\mathbb{Q}}=0\quad\text{for all }n>0.
\]
This is equivalent to torsion, not necessarily finiteness of the integral group or a single bounded exponent. The group $K_0(X)$ is excluded, and properness and smoothness are part of the chosen scope. Passing to rational coefficients is the vanishing formulation, not an assertion of integral vanishing.
\par\smallskip\textit{Formulation and concept references:} \href{https://www.sciencedirect.com/science/article/pii/S0022404923002001}{[S1]}.

\subsection{Short English statement}
Are all positive-degree algebraic K-groups of smooth projective varieties over finite fields torsion?
\subsection{Sources}
\begin{itemize}
\item[S1]On Parshin's conjecture.
\url{https://www.sciencedirect.com/science/article/pii/S0022404923002001}.
\textit{Formulation source.}
\end{itemize}
\end{document}
